\section{Fine-Tuning}
\label{sec:finetuning}

We fine-tune the base LLM on a domain corpus to give it fluency in SystemVerilog, SVA, and the assume-guarantee contract structure before the RL
stage (Section~\ref{sec:rl}).

\subsection{Data Sources}
\label{sec:ft-sources}

We compose the corpus from complementary sources.

\begin{table}[h]
\centering
\small
\begin{tabular}{llll}
\hline
Source & Purpose & Verified labels & Used for \\
\hline
HierSVA~\cite{hiersva2026} & Hierarchical contracts & Yes & SFT/RL \\
OpenTitan~\cite{opentitan} & Assume-guarantee structure & Yes & SFT/RL \\
HWSec-UNC~\cite{ahmad2024hwsec} & Security properties & Yes & SFT/RL \\
riscv-formal~\cite{riscvformal} & Compositional verification & Yes & SFT/RL \\
VERT~\cite{menon2025vert} & SVA syntax breadth & Mixed & SFT \\
Verixa~\cite{verixa2025} & Verification context & Mixed & SFT \\
NVIDIA CVDP~\cite{nvidiacvdp2025} & Specification and RTL & Mixed & SFT \\
VerilogEval~\cite{liu2023verilogeval} & English and RTL & No & SFT \\
\hline
\end{tabular}
\caption{Fine-tuning and reinforcement-learning data sources.}
\label{tab:ft-sources}
\end{table}

\subsection{Data augmentation}
\label{sec:ft-augmentation}

We expand each verified (specification, contract) pair along two axes:

\begin{itemize}
\item \textbf{Identifier renaming}: signal, port, and module names are substituted with sampled alternatives consistent with the design's naming conventions.
\item \textbf{Syntactic paraphrasing}: system verilog and SVA code are transformed into semantically equivalent variants using rule-based methods, such as reordering conjuncts, De Morgan's laws, and equivalent temporal operators.
\end{itemize}
